% Partie : théorie de la démonstration
% Chapitre : normalisation
% Section : segments et coupures

\documentclass[../../../include/open-logic-section]{subfiles}

\begin{document}

\olfileid[fr]{pt}{nor}{seg}

\olsection{Segments et coupures}

Une règle d'introduction suivie d'une règle d'élimination
constitue un détour dans une !!{proof} ; normaliser la preuve
consiste à supprimer ces détours. Les règles \Elim{\lor}
et \Elim{\lexists} compliquent toutefois la définition.
L'élimination ne suit pas nécessairement
\emph{immédiatement} l'introduction : une application de
\Elim{\lor} ou de \Elim{\lexists} peut les séparer.
Par exemple :
\[
\AxiomC{}
\DeduceC{$\lexists[x][!C(x)]$}
\AxiomC{$\Discharge{!A}{x}$}
\DeduceC{$!B$}
\DischargeRule{\Intro{\lif}}{x}
\UnaryInfC{$!A \lif !B$}
\DischargeRule{\Elim{\lexists}}{y}
\BinaryInfC{$!A \lif !B$}
\AxiomC{}
\DeduceC{$!A$}
\RightLabel{\Elim{\lif}}
\BinaryInfC{$!B$}
\DisplayProof
\]
Un nombre quelconque de ces règles peut même séparer
\Intro{\lif} de \Elim{\lif}, comme dans l'arbre suivant :
\[
\AxiomC{}
\DeduceC{$\lexists[x][!C(x)]$}
\AxiomC{}
\DeduceC{$!D \lor !E$}
\AxiomC{$\Discharge{!A}{x}$}
\DeduceC{$!B$}
\DischargeRule{\Intro{\lif}}{x}
\UnaryInfC{$!A \lif !B$}
\AxiomC{$\Discharge{!A}{x}$}
\DeduceC{$!B$}
\DischargeRule{\Intro{\lif}}{x}
\UnaryInfC{$!A \lif !B$}
\RightLabel{\Elim{\lor}}
\TrinaryInfC{$!A \lif !B$}
\DischargeRule{\Elim{\lexists}}{y}
\BinaryInfC{$!A \lif !B$}
\AxiomC{}
\DeduceC{$!A$}
\RightLabel{\Elim{\lif}}
\BinaryInfC{$!B$}
\DisplayProof
\]
Une même application de \Elim{\lif} peut donc participer
à plusieurs détours, puisqu'elle se trouve au-dessous de
plusieurs applications de \Intro{\lif}.
Pour décrire ces situations, nous considérons des suites
d'occurrences de !!{formula}s dans une preuve de \Log{N1i},
appelées \emph{segments}. Dans le deuxième arbre, les
occurrences de~$!A \lif !B$ se suivent d'une prémisse
mineure de $\Elim{\lor}$ ou de $\Elim{\lexists}$ à la
conclusion de cette inférence. Deux segments de ce type
comportent chacun trois occurrences de~$!A \lif !B$.

\begin{defn}
Un \emph{segment} d'une !!{proof}~$\delta$ de \Log{N1i}
est une suite non vide d'occurrences $!A_1$, \dots, $!A_n$
d'une même !!{formula}~$!A$ dans~$\delta$, telle que :
\begin{enumerate}
  \item $!A_1$ n'est la conclusion ni de \Elim{\lor}
  ni de \Elim{\lexists} ;
  \item $!A_n$ n'est une prémisse mineure ni de \Elim{\lor}
  ni de \Elim{\lexists} ;
  \item pour $i = 1$, \dots,~$n-1$, $!A_i$ est une
  prémisse mineure d'une application de \Elim{\lor}
  ou de \Elim{\lexists}, et $!A_{i+1}$ est la conclusion
  de cette même application.
\end{enumerate}
La \emph{longueur} du segment est~$n$.
\end{defn}

Tous les segments ne constituent pas des détours à
supprimer. Ceux qui en constituent sont appelés
\emph{coupures}.

\begin{defn}
Une \emph{coupure} est un segment dont la dernière
occurrence~$!A_n$ est la prémisse majeure d'une règle
d'élimination et qui satisfait l'une des deux conditions
suivantes : soit $n=1$ et $!A_1 = !A_n$ est la
conclusion d'une règle d'introduction ou de~\FalseInt,
soit~$n>1$.
\end{defn}

Un segment de coupure de longueur supérieure à un
ne commence donc pas nécessairement par la conclusion
d'une règle d'introduction. Pour une coupure de
longueur~$1$, en revanche, l'unique occurrence doit
être la conclusion d'une introduction \emph{ou
de~\FalseInt}, conformément à la définition.

\begin{defn}
Le \emph{rang} d'une coupure sur~$!A$ est~$\depth{!A}$.
Le \emph{rang de coupure}~$\cutrank{\delta}$ d'une
!!{proof}~$\delta$ est le rang maximal de ses coupures.
Une coupure est \emph{maximale} dans~$\delta$ si son
rang est~$\cutrank{\delta}$.
La \emph{longueur de coupure} de~$\delta$ est la somme
des longueurs de toutes ses coupures maximales.
Si la preuve ne comporte aucune coupure, nous
convenons que son rang et sa longueur de coupure
valent tous deux zéro.
\end{defn}













\end{document}
